\cotefile{18}
\scanpage{33}
Structure \inserted{de \unsure{vectoriel|naturel}} complexe sur $DR(M)$ ; mais on notera que
la bigraduation déduite par transport de structure
de celle de Hodge ne \uline{dépend pas de $u$}
(donc, si on remplace $u(\lambda)$ par $\overline{u(\lambda)}$, elle n'est \uline{pas}
remplacée par la symétrique !).

\inserted{Néanmoins, si on regarde,
pour un motif donné $M$,
le motif $\unsure{{}^{\sigma}M_{\underline{\mathbb C}}|{}^{\sigma}\!M_{\underline C}}$ ($\underline{\mathbb C}$ la clôt.
normale de $K$) alors $DR({}^{\sigma}M_{\underline{\mathbb C}})\xrightarrow{\sim}DR(M)_{\underline{\mathbb C}}$, les structures rationnelles \inserted{et \unsure{filtrations|filtration}} sont les \unsure{mêmes},
\unsure{Ainsi|Alors|mais} \deleted{la filtration} \unsure{composant|composé} avec $\sigma$, on trouve un \uline{$K$-isom}
$DR(M_{\underline{\mathbb C}})\simeq DR(M)$, conservant la $\unsure{Q|\mathbb Q}$-structure, et \uuline{renversant} filtration.}

\deleted{Supp.} Soit maintenant $K_0=\mathbb R$, et considérons les
motifs \inserted{$M$} \uline{sur} $\mathbb R$. Choisissons une clôture alg.
$K$ de $K_0=\mathbb R$. Alors $DR(M)_K=DR(M_K)$
est muni de \uline{deux façons} d'une $\mathbb R$-structure :
du fait qu'il est déduit de $DR(M)$ par $\mathbb R\unsure{-|\to}K$,
et puis : la construction donnée plus haut.
Soit $\sigma$ l'action de conjugaison complexe définie
plus haut, (N.B. $\sigma$ \uline{ne respecte pas} la \deleted{\unsure{bigr.}}\inserted{filtration})
$\tau$ celle définie par $DR(M)_K\simeq DR(M_K)$ (donc
\margin{$X(K)$}$\sigma$ respecte la filtration) -- $\sigma$ et $\tau$ sont
\uline{antilinéaires}. Alors $\sigma$ et $\tau$ commutent,
en d'autres termes $DR(M_K)^\sigma$ $\left(=H^*(X(K),\mathbb R)\ \text{si }M=\mathrm{Mot}^*(X)\right)$
est stable par $\tau$ : en effet, $\tau$ n'est autre
que l'action de $H^*(X(K),\mathbb R)$ définie par
l'action $\tau_X$ de $X(K)$ ! Alors $\sigma\tau=\tau\sigma=f^{\deleted{\unsure{*}}}_{\mathbb R}$
est l'action de $H^*(X(K),K)\simeq DR(\mathrm{Mot}^*(X))$
% archive sheet number: 32
\scanpage{34}
définie par l'\unsure{aut.|aut.t} \deleted{\unsure{de|$\tau_X$}} $\tau_X$ de $X(K)$). Ainsi, un
motif sur $\mathbb R$ définit :

a) Une structure de Hodge \inserted{$V$} sur $K$

b) Un autom. $f$\deleted{\unsure{pro}} \inserted{de carré 1} du $K$-espace vectoriel \deleted{sous-jacent}\inserted{$V_K$},
cet autom étant ``réel'' i.e. \deleted{\unsure{$f$|qui}} transformant $V_{\mathbb R}$ en $V_{\mathbb R}$,
\deleted{\unsure{On signifie}} et \unsure{$=$} \uline{rationnel} i.e. transforme $V$ en $V$

Ces deux données étant assujetties à la condition
\[
\left\{\begin{aligned}
 &( \deleted{f\unsure{\sigma}} )\,\sigma\ \text{inverse les filtrations naturelles}\\
 &\text{i.e. } f\sigma\ \text{échange la filtration donnée et l'opposée}
\end{aligned}\right.
\]

En termes de ``donnée de descente'', on peut \deleted{dire aussi}\inserted{exprimer}
aussi b) en disant qu'on a un isom
de la structure de Hodge $V$ avec la structure de
Hodge \deleted{\unsure{que}} ``conjuguée'' : c'est le frobénius réel.

Soit $G$ le groupe de Galois motivique de $M$ sur $\mathbb R$.
Alors $f$ correspond à un élément
\[
f_{\mathbb R}\in G(\mathbb Q)
\]
appelé \uline{frobénius réel}. C'est un élément d'ordre 2.
Son image dans $G/G^0$ engendre ce groupe. Cet
élément normalise $\underline i:S_{\mathbb R}\longrightarrow G_{\mathbb R}$, de façon précise
\[
f_{\mathbb R}\,i(\lambda)\,f_{\mathbb R}^{-1}=i(\deleted{\unsure{\sigma\lambda}}\lambda')
\qquad\text{où }\lambda\unsure{\sim|\mapsto}\lambda'\ \text{est la symétrie canonique de }S_{\mathbb R}
\]
donc sur $U_{\mathbb R}$ on a
\[
 f_{\mathbb R}\,i(\lambda)\,f_{\mathbb R}^{-1}=i(\lambda^{-1})
 \qquad(\lambda\in U_{\mathbb R}\deleted{\unsure{\mathbb R}}) \tag*{(*)}
\]

Le foncteur $M\longmapsto(\text{structure de Hodge associée, avec }f_{\mathbb R})$
est pleinement fidèle en vertu des conj. de Hodge.

% margin note written sideways
\margin{La structure en \deleted{\unsure{sont}} donne $G$ sur $\mathbb Q$, $\unsure{\varepsilon|E}:G\unsure{-|\to}G_m$, $i':S_{\mathbb R}\unsure{-|\to}G_{\mathbb R}$ (avec $i|G_{m\,\mathbb R}$ défini sur $\mathbb Q$ \unsure{etc|et}) $+\, f_{\mathbb R}\in G(\mathbb Q)$ satisfaisant la condition conj.}
% archive sheet number: 33
\scanpage{35}
% letterhead
\begin{center}
INSTITUT\\
DES HAUTES ÉTUDES SCIENTIFIQUES\\
BURES-SUR-YVETTE (S.-\&-O.)\\
928-48-53
\end{center}

Les \unsure{éléments|élémt} de $DR(M)$ sont les éléments de
$DR(M)_K\simeq V\otimes_{\mathbb Q}K$ de la forme $z=x+iy$ $(x,y\in V_{\mathbb R})$
avec $x=fx$, $y=-fy$ (on exprime $\tau z=z$ i.e. $f\sigma z=z$ i.e. $\sigma z=fz$)
\deleted{\unsure{C'est d.|Cela}} Donc l'espace $DR(M)$ admet une \inserted{$\mathbb Q$-structure de} façon naturelle,
\deleted{\unsure{de la structure}} on prend le sous-espace des
$x+iy$, $x,y\in V$, $x=fx$, $y=-fy$. Soit $W$ cet espace, il
mérite le nom de \uline{sous-espace de Betti \inserted{tordu}} de $DR(M)$.

\diagram{\begin{tikzcd}[ampersand replacement=\&,column sep=small,row sep=small]
H^*(X(K),\mathbb Q)=V \arrow[r,phantom,"\subset"] \arrow[d,phantom,"\cup"] \& DR(M)_K\simeq DR(M_K) \arrow[d,phantom,"\cup"]\\
H^*(X(\mathbb R),\mathbb Q)=V\cap W=V^f=W^{F} \arrow[r,phantom,"\subset"] \& W
\end{tikzcd}}

Il est stable par $f$. Mais
la formation de ce sous-espace
\unsure{ne|se} se comporte-t-il comme
il faut par $\otimes$ ? \uline{Oui}.

On a ainsi un deuxième foncteur fibre sur $M$
à valeurs sur $\mathbb Q$, mais il est isom \inserted{comme foncteur $\otimes$ \unsure{cohérent}} au
premier, \deleted{\unsure{une fois,}} comme on voit en écrivant
$V\simeq V^++V^-$, $W=W^++W^-$, et utilisant les
isomorphismes $V^+\simeq W^+$, $V^-\simeq W^-$ canoniques.
Cet isom $V\unsure{-|\to}W$ est \deleted{\unsure{dé.|f.}} induit par l'autom
de $DR(M_K)\simeq V_K\simeq W_K$ qui est l'identité sur
$V_K^+$, et la multiplication par $i$ sur $V_K^-$ ; Cependant, cet
isom \uline{n'est} \uline{pas} compatible avec le produit \uline{tensoriel}

Il semble que $\mathrm{Betti}\mathrm T$ et Betti ne sont \deleted{\unsure{pas}} peut être pas isom,
l'un est déduit \inserted{Betti \unsure{tordu}} de l'autre en tordant par le cocycle
défini par $f_{\mathbb R}$ ----
% archive sheet number: 34
\scanpage{36}
L'élément $\alpha\in P(\mathbb C)$ est relié à $f_\infty$ par
\[
\overline\alpha=f_\infty\alpha
\]
donc \deleted{$f_\infty$} à priori
\[
f_\infty=1\Longleftrightarrow\alpha\in P(\mathbb R)\quad\text{i.e.}\quad\mathbb R(\alpha)=\mathbb R
\]

\deleted{Comme on a la relation $f_\infty j_1 f_\infty^{-1}=j_2$, cette condition
implique $j_1=j_2$ donc $j_1=j_2=\frac12 i$, i.e. $\unsure{G_2|G}$ [si on
admet conj. de Hodge] $G^0\xrightarrow[\sim]{\varepsilon}G_m$}

Or on vérifie aussitôt
\[
\boxed{\varepsilon(f_\infty)=-1}
\]
ce qui nous montre qu'on a tjrs
\[
\boxed{f_\infty\ne1\quad\text{i.e.}\quad\mathbb R(\alpha)=\mathbb C}
\quad\text{(si $M$ contient le motif de Tate)}
\]

Si $M$ ne contient pas le motif de Tate, on voit
que $f_\infty=1\Longrightarrow j_1=j_2$ i.e. $j_1=j_2=\frac12 i$, donc en
vertu de la conj. de Hodge, il faut que
$G^0=\mathrm{Im}\,j_1$, donc ou bien $G^0=\{e\}$ \deleted{(\unsure{motifs de}
\unsure{dimension} $0$, alors $G=\{e\}$, \inserted{ici $G=G^0$, car} $G=\mathbb Z/2\mathbb Z$, dans ce
dernier cas $f_\infty$} \inserted{donc $G=1$}
[\uline{NB} (si $G$ est disconnexe, on sait que
$f_\infty\notin G^0(\mathbb Q)$)], ou bien $G^{\deleted{\unsure{0}}}=G_m/\mu_n$ \inserted{$n\geq1$},
$M$ est \deleted{\unsure{formé de}} engendré par $\mathbb Q(-n)$\} ; on trouve
en $=$ \unsure{temps} que $n$ est \uline{pair}. Donc on trouve

\diagram{\begin{tikzcd}[ampersand replacement=\&,column sep=small,row sep=small]
(f_\infty=1) \arrow[r,Leftrightarrow] \& (\mathbb R(\alpha)=\mathbb R) \arrow[r,Leftrightarrow] \&
\substack{M\ \text{est engendré par }\mathbb Q(-2n),\\\text{\unsure{pour|un} quelque }n\geq0.} \arrow[d,Leftrightarrow]\\
\& \& G\simeq G_m\ \uline{\text{ou}}\ G\simeq\{e\}
\end{tikzcd}}

% slanted
\margin{\uline{moyennant} \uline{conj. de Hodge}}

\uline{NB} $f_\infty$ ne peut être \uline{central} que si $G^0=\mathrm{Im}\,\deleted{\unsure{G|\varepsilon}}i$, ce qui
implique que $M$ est \deleted{engendré} contenu dans la
$\otimes$-cat. engendrée par $\mathbb Q$ tordu \uline{et} par $\mathbb Q(2)$.
Il revient au $\unsure{=|{-}}$ dire que $G$ est \uline{comm.}.
% archive sheet number: 35
\scanpage{37}
\uline{Motifs sur $\mathbb R$}

$\mathcal M$ \inserted{$\otimes$-}catégorie de motifs sur $\mathbb R$, $G=\underline{\mathrm{Aut}}_\otimes(T_B)$ groupe de Galois associé.

Notons qu'on a un $\otimes$-automorphisme d'ordre 2 de $T_B$,
qui sur la catégorie $\mathcal V_{\mathbb R}$ des var. proj. lisses $X$ sur $\mathbb R$
revient à l'automorphisme des $H^*(X(\mathbb C),\mathbb Q)$ défini par
l'autom $x\mapsto\bar x$ de $X(\mathbb C)$. D'où
\[
\boxed{f_{\deleted{\mathbb R}\infty}\in G(\mathbb Q)}\qquad(f_{\deleted{\mathbb R}\infty}^{2}=1),
\]
(``frobénius réel''). \deleted{On retrouve le sous-espace
$H^*_{DR}(X)$ de $H^*(X\unsure{,|_{\mathbb R},}\mathbb C)=1$} Si $\sigma_M$ est la conjugaison
complexe de $M\otimes_{\mathbb Q}\mathbb C$ [où $M\in\mathrm{ob}\,\mathrm{Rep}(C)$], alors
la conjugaison $\tau_M$ associée à la structure réelle
$T_{DR}(M)\subset M\otimes_{\mathbb Q}\mathbb C$ n'est autre que
\[
\tau_M=(f_{\deleted{\mathbb R}\infty})_M\,\sigma_M
=\sigma_M(f_{\deleted{\mathbb R}\infty})_M
\]

Donc on retrouve grâce à $f_{\deleted{\mathbb R}\infty}$ la structure mixte sur
les $M\otimes_{\mathbb R}\mathbb C$, donc aussi le torseur
\[
P=\underline{\mathrm{Isom}}_\otimes(T_B\otimes_{\mathbb Q}\mathbb R,T_{DR})
\]
sous $G_{\mathbb R}$ : c'est celui défini par le $1$-cocycle \deleted{de} non \unsure{abélien|normé|normalisé}
sur $\Gamma=\mathbb Z/2\mathbb Z$, à coeff. \deleted{\unsure{en|dans}} dans $G(\mathbb C)$, dont la valeur en
l'élément $\deleted{\ne}e$ est $f_{\deleted{\mathbb R}\infty}$. La forme tordue correspondante
de $G$ s'obtient en prenant sur $G_{\mathbb C}$ la donnée
de descente associée à \deleted{\unsure{$f$}} $\mathrm{int}(f_{\deleted{\mathbb R}\infty})\times_{\mathbb R}\sigma_{\mathbb C}$. \deleted{On doit
donc, pour} La relation entre $f_\infty$ et
\[
j:S_{\mathbb R}\longrightarrow G_{\mathbb R}
\]
est que
\[
\boxed{f_\infty j(\lambda) f_\infty^{-1}=j(\lambda^{\deleted{\unsure{\sigma}}\,\prime})},
\]
où $\lambda\mapsto\lambda'$ est l'autom can. de $S_{\mathbb R}$,
(donnant $\lambda\mapsto\bar\lambda$ sur $S_{\mathbb R}(\mathbb R)\simeq\mathbb C^*$)
% archive sheet number: 36
\scanpage{38}
D'autre part, on aura ici
\[
\beta\in\unsure{Q|\mathbb Q}(\mathbb R)\subset P'(\mathbb R)
\]
donc
\[
\mathbb R(\gamma)=\mathbb R(\beta,\gamma)=\mathbb R(\beta,\alpha)=\mathbb R(\alpha)
\]
et par suite
\[
\mathbb R(\gamma)=\mathbb R(\alpha)
\]
qui est donc $=\mathbb C$ sauf dans le cas exceptionnel précisé
plus haut.

\uline{Remarque} Soit $M$ un motif sur $\mathbb R$, correspondant à un \deleted{\unsure{$G$-module}
\unsure{aussi noté}} $T_B(M)$. Alors on ne peut avoir $f_{\infty M}=\mathrm{id}_M$ que
si $M$ est une somme de puissances tensorielles
de $\mathbb Q(2)$.

\deleted{La relation
\[
 f_\infty j(\lambda) f_\infty^{-1}=j(\lambda')
\]
implique, faisant $\lambda=C\in\unsure{S|D}(\mathbb R$
\[
f_\infty C f_\infty^{-1}=C^{-1}=\unsure{\eta|y}C\qquad(\eta=i(-1))
\]
donc
\[
 f_\infty\in\underline{\mathrm{Norm}}_{G_{\mathbb R}}(C,C^{-1})
\]
\deleted{\unsure{Or, on a} $K'=\underline{\mathrm{Norm}}_{G_{\mathbb R}}(C,C^{-1})\supset\underline{\mathrm{Norm}}_{G_{\mathbb R}}(C)=K$ qui est un
sous groupe d'indice 1 ou 2. On sait que $K(\mathbb R)$ est un
compact maximal dans $G(\mathbb R)$} Considérons $G/i(G_m)=\dot G$
$[=\mathrm{Ker}\,\varepsilon/i(\mu_2)]$, et les images \deleted{\unsure{de}} $\dot f_\infty$, $\dot C$ de $f_\infty$, $C$,
donc on a
\[
 \dot f_\infty\in\underline{\mathrm{Norm}}_{\dot G_{\mathbb R}}(\dot C,\dot C^{-1})\overset{\mathrm{dfn}}=K'
\]
Or \deleted{on sait \unsure{qu'on}} sait
\[
 K=\underline{\mathrm{Norm}}_{\dot G_{\mathbb R}}(\dot C)\subset K'
\]}
% archive sheet number: 37
\scanpage{39}
\deleted{\unsure{Soit}}

Soit $K$ un corps topologique isomorphe à $\mathbb C$, $X$ un
schéma algébrique défini sur $K$. Alors $X_{\mathrm{top}}$ est défini,
et aussi $H^*(X_{\mathrm{top}},K)$. D'autre part, \inserted{si $X$ est non sing.,} la cohomologie
de Hodge est définie, soit $H^*(X)$ [indépendant de
la structure topologique de $K$]. Alors le choix
d'un isom top. $u:K\unsure{-|\to}\mathbb C$ définit un isomorphisme
\[
H^*(X_{\mathrm{top}},K)\xrightarrow{\sim} H^*(X)
\]
et on vérifie aussitôt que cet isomorphisme ne
dépend pas du choix de $u$, [En d'autres termes,
si $K=\mathbb C$, alors les deux homomorphismes
\diagram{\begin{tikzcd}[ampersand replacement=\&,row sep=small]
\& H^*(X) \arrow[dd,"\simeq\ (\text{antilinéaire})"]\\
H^*(X_{\mathrm{top}},\mathbb C) \arrow[ur] \arrow[dr] \arrow[r,phantom,"\text{non comm.}" description] \& \\
\& H^*(X^c)
\end{tikzcd}}
(où $c:\mathbb C-\mathbb C$ est la conjugaison complexe)
sont \deleted{conjugués} déduits
l'un de l'autre par application de $c$ aux coefficients de $H^*(X_{\mathrm{top}},\mathbb C)$]

\inserted{Cela nous permet, notant qu'on a un \unsure{isom|hom|homom.} injectif continu unique
$\unsure{u|i}:\mathbb R\to\mathbb C$, de définir $H^*(X_{\mathrm{top}},\mathbb R)\hookrightarrow H^*(X)$ (même si $X$ pas définissable
% margin note written sideways
sur $\mathbb R$ !) d'où un sous-\inserted{$\mathbb R$-}espace
de $H^*(X)$, \deleted{défini sur $\mathbb R$, \unsure{et}} noté $H^*(X_{\unsure{\tau|+|?}},\mathbb R)$
\deleted{$\unsure{H^*(X)^+}$}, ou \deleted{\unsure{cet} \unsure{dépend de la topologie de} $K$},
\unsure{aut.nt|anti} $H^*(X)^-=\sqrt{-1}H^*(X)^+$.}

Supposons $X$ projective non singulière, alors
le choix d'un isom top. $u:K\xrightarrow{\sim}\mathbb C$ définit
une décomposition
\[
\varphi:\coprod_{p,q}H^q(X,\Omega_X^p)\xrightarrow{\sim}H^*(X)
\]
compatible avec la filtration de $H^*(X)$. On
vérifie encore que $\varphi$ est indépendant du choix de $u$.
% archive sheet number: 38
\scanpage{40}
Enfin, supposons maintenant que $K$ soit isomorphe
$\mathbb R$ [on peut supposer $K=\mathbb R$, car $\mathbb R$ n'a pas d'
autom. non trivial] soit $K'$ une clôture
algébrique de $K$, $\pi=\mathbb Z/2\mathbb Z$ son groupe de Galois,
$X$ lisse \deleted{\unsure{algébrique|projectif}} [projectif] sur $K$. Alors on peut appliquer
ce qui précède à $X_{K'}$ sur $K'$,
\deleted{$H^*(X_{K'\,\mathrm{top}},\mathbb R)\unsure{\xrightarrow{\sim}|\longrightarrow} H^*(X)\otimes_K K'$}, on trouve
\[
\left[\varphi_{K'}:\coprod_{p,q}H^q(X,\Omega_X^p)\otimes_K K'
 \xrightarrow{\sim}H^*(X_{\deleted{\unsure{K'}}})\otimes_K K'\right]
\]
et d'après ce qui a été vu plus haut,
\unsure{cet|ces} isomorphisme\deleted{\unsure{s}} \uline{commute\deleted{\unsure{nt}} à} $\pi$, donc provient
d'un isomorphisme
\[
 \varphi_1:\coprod_{p,q} H^q(X,\Omega_X^p)\xrightarrow{\sim}H^*(X)
\]

% margin note written sideways
\margin{On a aussi un $\mathbb R$-isom canonique $(\pi=\mathbb Z/2\mathbb Z)$
$H^*(X_{\mathbb C\,\mathrm{top}},\mathbb R)\otimes_{\mathbb R}\mathbb C\xrightarrow{\sim}H^*(X)\otimes_{\mathbb R}\mathbb C$,
où $\pi$ opère à gauche via son opération sur $X_{\mathbb C\,\mathrm{top}}$,
\uline{et} sur $\mathbb C$ dans $\otimes_{\mathbb R}\mathbb C$
d'où en prenant les invariants sous $\pi$ une décomposition
canonique $H^*(X)\simeq H^*(X)^+\oplus H^*(X)^-$,
\deleted{$\unsure{H^*(X,\mathbb R)}$}
$H^*(X)^+\simeq H^*(X_{\mathbb C\,\mathrm{top}},\mathbb R)^\pi\simeq H^*(X_{\mathrm{top}}=X_{\mathbb C\,\mathrm{top}}/\pi_0\,;\mathbb R)$,
$H^*(X)^-=H^*(X_{\mathbb C\,\mathrm{top}},\mathbb R)^-$.}

En résumé, chaque fois que $X$ est
projectif non singulier sur un corps
\deleted{\unsure{top. ou|top.}} top. loc. \unsure{cpt|cpt.} complet archimédien,
on a une décomposition $\varphi$ comme
dessus de $H^*(X)$, compatible avec la
filtration. \deleted{Si} On pose
\[
H^{p,q}(X)=\varphi\big(H^q(X,\Omega_X^p)\big)\subset \deleted{\unsure{F|E}}F^p H^q(X)
\]

Si $K$ n'est pas topologique, alors pour
tout choix d'une topologie \unsure{forterie|forte|localement} archimédienne
$T$ sur $K$, on déduit sur
% archive sheet number: 39
\scanpage{41}
\[
H^*(X)_T=H^*(X)\otimes_K\widehat K_T.
\]
une décomposition,
\[
\boxed{\varphi_T:\coprod_{p,q}H^q(X,\Omega_X^p)_T\longrightarrow H^*(X)_T}
\]

Cela permet de définir
\[
H_T^{p,q}(X)=\varphi_T\big(H^q(X,\Omega_X^p)_T\big)\cap H^*(X)
 \subset F^p\big(H^q(X)\big)
\]
qui est un \uline{sous}-\inserted{$K$-}\uline{espace vectoriel de} $H^*(X)$,
défini canoniquement par la donnée de $T$. Nous
\unsure{poserons|posons|pouvons} de $=$
\[
H_T^{pq}(X,\mathbb R)=\varphi_T\big(H^q(X,\Omega_X^p)_T\big)
\cap H^*(X_{\widehat K_T},\mathbb R)\cap H^*(X)
\]
\deleted{c'est un \unsure{sous} \unsure{$K$} \unsure{$\mathbb R\cap K_T$}-espace vectoriel de}
c'est encore un sous-\inserted{$K$-}espace vectoriel de $H^*(X)$
si $\widehat K_T\simeq\mathbb R$ ; si\deleted{non} \inserted{$\widehat K_T\simeq\mathbb C$} c'est un sous-\deleted{\unsure{$R_{\mathbb R}$|$\mathbb R_T$}}$R_T$-espace vectoriel, où
\[
R_T=\mathbb R\cap K\qquad\text{dans }\widehat K_T\simeq\mathbb C
\]
[sous-\deleted{\unsure{espace}}\inserted{corps} \uline{réel} pour $T$] ; à fortiori, \deleted{\unsure{c'est}} dans tous les
cas, on obtient un $K_r^{\unsure{*}}$-espace \unsure{vectoriel|vectorel}, où
\[
K_r\simeq\bigcap_T R_T\simeq\text{plus \unsure{grand|gd} sous-corps tot. réel de }K
\]
(NB $K_r$ est une ext. algébrique du corps premier $\mathbb Q$).
% archive sheet number: 40
\scanpage{42}
\unsure{Variations|Variation} du réseau discret en \unsure{terme|termes} d'une place archimédienne,

Critère de transcendance
% archive sheet number: 41
\scanpage{43}
\uline{Pbs}

\deleted{\unsure{Pbs}}

\diagram{\(\begin{array}{l@{\quad}l}
\text{1) Betti, De Rham} &
\left\{\begin{array}{ll}
\text{a)} & \text{Comparaison ``réseaux'' de De Rham et \deleted{\unsure{avec|de}} de Betti}\\
\text{b)} & \text{Comparaison réseaux de Betti pour topologies différentes?}\\
\text{c)} & \text{Comparaison structures réelles, \uline{resp.} \uline{splittings}, \inserted{(dans De Rham)} pour top. diff.}
\end{array}\right.
\end{array}\)}

2) Betti, $l$-adique \deleted{Comp} Comparaison de réseaux de Betti pour top. diff.
(dans \uline{ch. $l$-adique}).

3) DeRham, $p$-adique \quad Sur un corps valué \inserted{(complet d'\unsure{inégales|inégale} car.)} non archimédien\inserted{, corps résiduel car $p$},
comparaison possible entre cohomologie $p$-adique et
cohomologie de De Rham ?

4) Définition de cohomologie $p$-adique en car $p>0$ ?
% archive sheet number: 42
\scanpage{44}
\mspage{32}

the induction assumption that the right hand side is a polynomial in $n$, of degree not higher than $s-1$. Therefore, $\unsure{X|\chi}(V^*,K_n^*)$ is a polynomial in $n$, of degree not higher than $s$.

The line from above will mean that we have a direct image, corresponding to the mapping $f$. There occurs an imbedding $K\to\overline K'$ [See $({}^{8})$]. To the imbedding $\unsure{O|0}'_V\to O_{V*}$ corresponds the homomorphism
\[
\handwritten{K'\otimes_{O'_V}O'_V\longrightarrow K'\otimes_{O'_V}O_{V^{*}}}
\]
or, what comes to the same thing, $K'\to K^*$.

Turning to the direct images, we obtain the homomorphism $\overline K\to\overline K^*$ and hence the homomorphism $\overline K\to\overline K{*}$.

Let $W$ be the variety of singularities of the variety $V$,
\[
\handwritten{U_W=V-W,\quad U'_W=f^{-1}(U_W).}
\]
$f$ establishes biregular correspondence between $U_W$ and $U'_W$.

\margin{p. 192}
Evidently, the homomorphisms $K|U_W\to\overline K'|U_W$ and $O'_V|_{U'_W}\to O_{V^{*}}|_{U'_W}$ are isomorphisms. Hence it follows that $K'|_{U'_W}\to(K'\handwritten{\otimes} O_{V^{*}})_{U'_W}$ are also isomorphisms.
% archive sheet number: 43
\scanpage{45}
$C,C'$ courbes elliptiques sans mult. complexe, avec $\unsure{\tau,\tau'|j,j'}$ \unsure{algébriques|algébriq.}

\[
\begin{aligned}
&C\times C'=A\\
&G\overset{\inserted{\text{\unsure{isogène|isog.|isogénie}}}}{\deleted{\simeq}}\mathrm{Sl}(2)\times\mathrm{Sl}(2)\times G_m\qquad\text{de dim }7\\
&B\simeq B(2)\times B(2)\times G_m\qquad\text{de dim }5
\end{aligned}
\]

Alors $K$ de degré de transcendance $\leq5$ sur $k$
donc de degré de transcendance $\leq5$ sur $k(\pi)$
[donc pas égal à $\dim G-1=6$].

\uline{NB} Cela joue pour \uline{\unsure{tte|toute|ttes}} immersion $k\subset\mathbb C$ !

% margin note written upside down
\deleted{On définit aussi
\[
\check H^q(\unsure{\widehat E|\widehat{\mathcal E}}/X,\unsure{\widetilde F|F})
\simeq\deleted{\text{\unsure{donne à}}}\check H^0(\widehat C_{/X},F)
\left[=\check{\mathcal H}^0(F)(X)\quad\text{si }X\in\mathrm{ob}\,C\right]
\]
\inserted{\uline{préfaisceaux}}
\unsure{Évident}
\[
 \check{\mathcal H}^q(F)(X)=\check H^q(\widehat C_{/X},F)
\]
\inserted{\uuline{préfaisceau}}}
% archive sheet number: 44
\scanpage{46}
\mspage{- 18 -}
n\textsuperscript{o} 370

\uline{positive sur} $\handwritten{\mathbb{Q}}_p$ \uline{telle que toute classe suivant} $p^n\handwritten{\mathbb{Z}}_p$ ($n$ \uline{entier} $\handwritten{\geq}0$) \uline{soit de mesure} $p^{-n}$.

La restriction de $\handwritten{\alpha}$ à $\handwritten{\mathbb{Z}}_p$ est évidemment une mesure de Haar sur $\handwritten{\mathbb{Z}}_p$.

6. \uline{Définition locale d'une mesure de Haar}.

PROPOSITION 7.- \uline{Soient} $G$ \uline{un groupe localement compact}, $V$ \uline{une partie ouverte de} $G$,
$\handwritten{\mu}$ \uline{une mesure positive non nulle sur} $V$, \uline{ayant la propriété suivante} : \uline{si} $U$ \uline{est une
partie \unsure{ouverte} de} $V$ \uline{et si} $s\handwritten{\in} G$ \uline{est tel que} $sU\handwritten{\subset} V$, \uline{l'image de la mesure} $\handwritten{\mu}_U$,
\uline{induite par} $\handwritten{\mu}$ \uline{sur} $U$, \uline{par l'homéomorphisme} $\unsure{x|\pi}\to sx$ \uline{de} $U$ \uline{sur} $sU$, \uline{est} $\handwritten{\mu}_{sU}$. \uline{Alors il
existe sur} $G$ \uline{une mesure de Haar à gauche} $\handwritten{\alpha}$ \uline{et une seule induisant} $\handwritten{\mu}$ \uline{sur} $V$.

Pour tout $s\handwritten{\in} G$, soit $\handwritten{\mu}_s$ l'image de $\handwritten{\mu}$ par l'homéomorphisme $x\to sx$ de $V$ sur
$sV$. La restriction de $\handwritten{\mu}_s$ à $V\handwritten{\cap} sV$ est l'image de $\handwritten{\mu}_{s^{-1}(V)\handwritten{\cap} V}$ par la restriction de
$x\to sx$ à $s^{-1}(V)\handwritten{\cap} V$ ; par hypothèse, cette image est $\handwritten{\mu}_{V\handwritten{\cap} sV}$. Par
translation, on en conclut que $\handwritten{\mu}_s$ et $\handwritten{\mu}_t$ ont même restriction à $sV\handwritten{\cap} tV$ quels que
soient $s$, $t$. D'après la prop.1 du chap. III, \S~3, n\textsuperscript{o} 1, il existe donc une mesure $\handwritten{\alpha}$
sur $G$ induisant $\handwritten{\mu}_s$ sur $sV$ quel que soit $s$. Il est clair que $\handwritten{\alpha}$ est l'unique mesure
de Haar à gauche sur $G$ induisant $\handwritten{\mu}$ sur $V$.

COROLLAIRE.- \uline{Soient} $G$, $G'$ \uline{deux groupes localement compacts}, $V$ (resp. $V'$) \uline{un voisinage ouvert de l'élément neutre de} $G$ (resp. $G'$), $\handwritten{\varphi}$ \uline{un isomorphisme local de} $G'$ \uline{à} $G$,
\uline{défini sur} $V'$, \uline{tel que} $\handwritten{\varphi}(V')=V$. \uline{Soient} $\handwritten{\alpha}'$ \uline{une mesure de Haar à gauche sur} $G'$,
\uline{et} $\handwritten{\alpha}'_{V'}$ \uline{sa restriction à} $V'$. \uline{Alors} $\handwritten{\varphi}(\handwritten{\alpha}'_{V'})$ \uline{est la restriction à} $V$ \uline{d'une mesure de
Haar à gauche unique} $\handwritten{\alpha}$ \uline{sur} $G$.

Soit $V_1$ un voisinage ouvert de $e$ dans $G$ tel que $V_1V_1^{-1}\handwritten{\subset} V$. Soit $\handwritten{\mu}$ la restriction de $\handwritten{\varphi}(\handwritten{\alpha}'_{V'})$ à $V_1$. Soient $U$ une partie ouverte de $V_1$ et $s\handwritten{\in} G$ tels que $sU\handwritten{\subset} V_1$.
On a $s\handwritten{\in} V_1V_1^{-1}\handwritten{\subset} V$, donc $s=\handwritten{\varphi}(s')$ avec un $s'\handwritten{\in} V'$. Soit $x\handwritten{\in} U$. On a $x=\handwritten{\varphi}(x')$
avec un $x'\handwritten{\in} V'$, donc $sx=\handwritten{\varphi}(s')\handwritten{\varphi}(x')=\handwritten{\varphi}(s'x')$ puisque $sx\handwritten{\in} sU\handwritten{\subset} V$. Comme les
translations à gauche dans $G'$ conservent $\handwritten{\alpha}'$, on voit que $V_1$ et $\handwritten{\mu}$ satisfont aux
conditions de la prop. 7. Soit $\handwritten{\alpha}$ la mesure de Haar à gauche sur $G$ induisant $\handwritten{\mu}$ sur
% archive sheet number: 45
\scanpage{47}
$M$ motif sur $k$ \inserted{alg. clos}, $k$ isom à un sous-corps de $\mathbb C$

Pour toute place $v$ de $k$, on a $M_v^{\mathbb Z}$ qui est
un sous-réseau de $\prod_{l\ \mathrm{fini}}M(l)$. On veut établir
les relations entre les $M_v^{\mathbb Z}$ pour diverses
valeurs $v_1\ldots v_n$ de $v$. Choisissons un
bon foncteur de référence (p. ex. défini par $v_0$)
$\unsure{\xi|\zeta}\ M\unsure{\sim|\mapsto|\rightsquigarrow} M_{\xi_0}\quad(\xi_0\in T(\mathbb Q))$ avec \deleted{des} isom \unsure{fonct.|fonctoriels|foncto}
$\xi_{0\deleted{\unsure{\mathbb Q_l}}}(\deleted{\unsure{H}}M)_{\mathbb Q_l}\simeq M_l\otimes_{\mathbb Z_l}\mathbb Q_l$. Alors $v_1,\ldots,v_n$ définissent des
fibrés principaux homogènes
\[
P_i=\underline{\mathrm{Isom}}(\xi,\xi_{v_i})
\]
sous $G_{M\xi_0}^{\unsure{k|f|\xi}}$ (sur $\mathbb Q$) et on a un
point de $(\prod P_i)(\unsure{\widehat{\overline{\mathbb Q}}|\widehat{\mathbb Q}})$, où $\widehat{\mathbb Q}=\prod_{l\ \mathrm{fini}}\mathbb Q_l$,
d'où un pt de $[(\prod P_i)/G](\widehat{\mathbb Q})$.
% slanted
\margin{$G$ opère diagonalement}

On peut aussi considérer que l'on a un
point canonique dans \deleted{$\mathrm{Gl}(n)\times\ldots\times\mathrm{Gl}(n)$}
\deleted{\unsure{$\mathrm{Gl}$}}
\[
\begin{aligned}
&\deleted{\backslash[\mathrm{Gl}(n)\times\ldots\times\mathrm{Gl}(n)](\widehat{\mathbb Q})/\mathrm{Gl}(n)(\widehat{\mathbb Q})}\\
&\deleted{(\mathrm{Gl}(n)\times\ldots\times\mathrm{Gl}(n))/\mathrm{Gl}(n)}\\
&\mathrm{Gl}(m,\mathbb Z)\times\ldots\mathrm{Gl}(m,\mathbb Z)\backslash
\mathrm{Gl}(m,\widehat{\mathbb Q})\times\ldots\times\mathrm{Gl}(m,\widehat{\mathbb Q})/
\mathrm{Gl}(m,\widehat{\mathbb Q})\qquad m=\text{rang }M
\end{aligned}
\]

Les considérations qui précèdent montrent que
le degré de transcendance du \unsure{pt|pts} en question
% archive sheet number: 46
